\documentclass[11pt]{article}
\usepackage[T1]{fontenc}
\usepackage[utf8]{inputenc}
\usepackage{lmodern}
\usepackage[a4paper,margin=1.1in]{geometry}
\usepackage{microtype}
\usepackage[hidelinks]{hyperref}
\usepackage{parskip}
\setlength{\parskip}{0.7em}

\title{\textbf{My take on the direction suggested by the new AMIE study in \textit{The Lancet}}}
\author{Spencer Dorn}
\date{Posted on LinkedIn, 13 October 2026\\[0.3em]
\small\href{https://www.linkedin.com/posts/spencerdorn_my-take-on-the-direction-suggested-by-the-share-7515908032132751361-FaUi/}{linkedin.com/posts/spencerdorn\_my-take-on-the-direction-suggested-by-the-share-7515908032132751361-FaUi}}

\begin{document}
\maketitle

My take on the direction suggested by the new AMIE study published in \textit{The Lancet}.

In brief, nearly 100 patients first interacted with AMIE, Google's conversational AI system, which generated a differential diagnosis and management plan. The patients then saw a real clinician who had access to AMIE's output. The researchers followed them to determine the final diagnosis.

There were no safety concerns, AMIE's differential included the final diagnosis 90\% of the time, and 75\% of PCPs found its output helpful.

By now, it should be quite clear that we're quickly moving beyond using AI to automate the notes we write about doctor--patient conversations to automating parts of the conversation itself. And from passive decision support toward something much more active.

While AMIE isn't ready for prime time, companies like K Health and Akido Care offer early signals of what this might look like in practice. I've written about both. Perhaps the cleanest analogy is how teaching attendings like me have long worked with three or four trainees in clinic at once.

That allows us to care for far more patients than we could alone. And done well, the additional set of eyes and ears can make us better.

But our patients may not feel they know us as well, and we don't know them as well either. We have to trust the person---or, increasingly, the machine---doing some of the work for us, even though we remain ultimately responsible.

We shouldn't romanticize the status quo. With demand far outstripping capacity, we need innovations like this. In my own clinic, I spend a full hour with new patients, but my practice is deeply specialized, and new patients wait many months to see me.

But we'll also need to be deliberate about what we give up. I was fortunate to recently hear Abraham Verghese talk about the ``iPatient'' and ``iDoctor.'' Increasingly, our digital representations will interact with one another, almost like shadows on the wall of Plato's cave.

There's enormous potential here. But it won't all be good, at least not for everyone.

\vfill
\begin{center}\footnotesize
Text reproduced from the original post on LinkedIn for easier reading. All words are Spencer Dorn's.
\end{center}
\end{document}
